Los datos se portaron como esperábamos. Mientras el punto de contacto se acerca a P4 la longitud del camino disminuye, y en cuanto lo rebasa vuelve a crecer, que es justo la forma de la función de la ecuación \eqref{eq:longitud} y la razón por la que la gráfica parece un valle con un fondo bien definido.

Ahora, ¿qué tan cerca quedó el mínimo medido del calculado?

Bastante. La parábola pone el mínimo en 12.63 cm y nosotros lo encontramos en 12.2 cm, así que el error porcentual es
\begin{equation}
E = \frac{\left| 12.63 - 12.2 \right|}{12.63} \times 100 = 3.4 \text{ por ciento}.
\label{eq:error}
\end{equation}
En la longitud la coincidencia es todavía mejor, porque 39.4 cm contra 39.38 cm es una diferencia de apenas dos décimas de milímetro. Esto tiene sentido. Cerca del mínimo la curva es casi plana, de modo que mover el punto unos milímetros a la izquierda o a la derecha cambia muy poco la longitud, y eso hace que la longitud mínima sea fácil de acertar pero su posición exacta no tanto.

Las diferencias que sí aparecieron vienen, sobre todo, de la propia técnica. La regla solo marca milímetros, el grosor del lápiz agrega incertidumbre al ubicar cada punto, la imagen A' se construyó a mano y cualquier pequeña inclinación en la perpendicular la desplaza, y además solo teníamos siete puntos, varios de ellos lejos del mínimo; con más puntos alrededor de P4 la parábola habría quedado mejor definida.

También hay que aclarar algo. Una parábola no es la función real, que es la suma de dos raíces cuadradas de la ecuación \eqref{eq:longitud}, sino una aproximación que funciona bien cerca del mínimo, y lejos de él, en P7 por ejemplo, ya se separa un poco de los datos.

En conclusión, comprobamos que de todos los caminos posibles entre A y B pasando por el espejo, el más corto es el que corresponde a la reflexión real, y que en ese camino A', P y B quedan alineados, lo cual obliga a que el ángulo de incidencia sea igual al de reflexión. Es decir, el principio de Fermat se cumple.

Fue una práctica sencilla, pero deja claro algo importante: la ley de la reflexión no es una regla aislada que haya que memorizar, sino la consecuencia natural de que la luz busque el camino más corto.
